Diagramy zobrazují ztrátové číslo $\varepsilon''$ proti reálné složce permitivity $\varepsilon'$; frekvence je parametr.
\begin{itemize}
\item \textbf{Cole--Cole:} kruhový oblouk odpovídající symetrickému rozdělení relaxačních dob; při rozšíření rozdělení leží střed kružnice pod reálnou osou. Debyeův případ s jedinou relaxační dobou je půlkružnice se středem na ose.
\item \textbf{Havriliak--Negami:} umožňuje kromě rozšíření i asymetrii rozdělení relaxačních dob, tedy nesymetrický oblouk.
\end{itemize}
\textbf{Havriliakův--Negamiho model lépe vystihuje reálné dielektrikum}, protože jeho polarizující částice nemají jedinou relaxační dobu a jejich rozdělení může být nesymetrické.
